
\documentclass[12pt]{article}
\usepackage{geometry}

\geometry{paperwidth=18cm, paperheight=8.7cm, margin=0.2cm}

\usepackage{bfcours}
\usepackage{bfcours-fonts}
\usepackage{pagecolor} % Pour colorer le fond de la page
%\usetikzlibrary{backgrounds}
\usepackage{expl3} % Nécessaire pour utiliser le langage LaTeX3

% Retourne la valeur absolue d'un entier
\newcommand{\absnum}[1]{%
    \ifnum #1<0 \number\numexpr -#1\relax\else \number\numexpr #1\relax\fi%
}

% Retourne le minimum de deux entiers
\newcommand{\minnum}[2]{%
    \ifnum #1<#2 #1\else #2\fi%
}

% Calcul récursif du PGCD de deux entiers (les arguments doivent être entiers)
\newcommand{\gcdcalc}[2]{%
    \ifnum #2=0 %
    \absnum{#1}%
    \else %
    \gcdcalc{#2}{\the\numexpr #1 - (#1/#2)*#2 \relax}%
    \fi%
}
\ExplSyntaxOn
% Définition de la commande \mprefix
\cs_new:Npn \mprefix #1
{
    \int_case:nnF {#1}
    {
        {-12} { \text{pico} }
        {-9}  { \text{nano} }
        {-6}  { \text{micro} }
        {-3}  { \text{milli} }
        {-2}  { \text{centi} }
        {-1}  { \text{deci} }
        {0}   { \text{} } % Aucun préfixe pour la puissance 0
        {1}   { \text{deca} }
        {2}   { \text{hecto} }
        {3}   { \text{kilo} }
        {6}   { \text{mega} }
        {9}   { \text{giga} }
        {12}  { \text{tera} }
    }
    { \text{Valeur inconnue} } % Gestion des cas non pris en charge
}
\ExplSyntaxOff


% Calcul et affichage des diviseurs avec Lua
\begin{luacode*}
    function generate_divisors_table(n)
    local seen = {} -- Table pour suivre les nombres déjà rencontrés
    tex.print("\\noindent\\begin{tabular}{c|c}")
    for i = 1, n do
    if n % i == 0 then
    local complement = math.floor(n / i) -- Calcul de la partie entière
    -- Vérifier si le complément est déjà dans la table 'seen'
    if seen[complement] then break end
    -- Ajouter le diviseur et le complément dans la table
    seen[i] = true
    tex.print(i .. " & " .. complement .. " \\\\")
    end
    end
    tex.print("\\end{tabular}")
    end
\end{luacode*}

\newcommand{\divisorstable}[1]{%
    \directlua{generate_divisors_table(#1)}%
}

% Commande pour s'adapter en fonction du choix
\newcommand{\marchandcmd}[1]{%
    \IfStrEqCase{#1}{%
        {marchand}{\xdef\munite{articles}\def\intro{}\def\vend{vend}\def\tousSes{$\proportion$ de ses }\def\oeufs{$\quantite$\ articles}\def\labelvente{le nombre d'articles vendus}}
        {agriculteur}{\xdef\munite{tonnes}\def\intro{Suite à des intempéries,}\def\vend{perd}\def\tousSes{les $\proportion$èmes}\def\oeufs{sa récolte, estimée à \quantite\  tonnes}\def\labelvente{la masse perdue}}
        {boulanger}{\xdef\munite{baguettes}\def\intro{}\def\vend{vend}\def\tousSes{les $\proportion$èmes}\def\oeufs{de ses $\quantite$ baguettes}\def\labelvente{le nombre de baguettes vendues}}
    }[\textcolor{red}{\textbf{Métier inconnu : #1}}]% Si aucune correspondance
}

% Définition de la commande criteredivisibilite
\newcommand{\criteredivisibilite}[2]{%
    % #1 est le diviseur (n), #2 est le nombre à tester (N)
    \def\diviseur{#1}%
    \def\nombre{#2}%
    
    % Calcul du dernier chiffre du nombre
    \def\dernierchiffre{\intcalcMod{#2}{10}}
    
    % Calcul des deux derniers chiffres du nombre
    \def\deuxdernierchiffres{\intcalcMod{#2}{100}}
    
    % Vérifier quel diviseur a été spécifié
    \ifthenelse{\equal{\diviseur}{2}}{%
        % Cas 2
        \textbf{Critère de divisibilité par 2 :} Un nombre est divisible par 2 si son dernier chiffre est pair (0, 2, 4, 6 ou 8).
        
        \textbf{Application à \nombre :} Le dernier chiffre de \nombre{} est \dernierchiffre.
        
        \ifthenelse{\equal{\intcalcMod{\nombre}{2}}{0}}{%
            Comme le dernier chiffre est \dernierchiffre{} (qui est pair), \nombre{} est divisible par 2.
            }{%
            Comme le dernier chiffre est \dernierchiffre{} (qui n'est pas pair), \nombre{} n'est pas divisible par 2.
        }%
        }{%
        \ifthenelse{\equal{\diviseur}{3}}{%
            % Cas 3
            \textbf{Critère de divisibilité par 3 :} Un nombre est divisible par 3 si la somme de ses chiffres est divisible par 3.
            
            \textbf{Application à \nombre :}
            
            Pour vérifier si \nombre{} est divisible par 3, calculons la somme de ses chiffres.
            
            % Dans un cas réel, on calculerait la somme des chiffres
            \ifthenelse{\equal{\intcalcMod{\nombre}{3}}{0}}{%
                La somme des chiffres de \nombre{} est divisible par 3, donc \nombre{} est divisible par 3.
                }{%
                La somme des chiffres de \nombre{} n'est pas divisible par 3, donc \nombre{} n'est pas divisible par 3.
            }%
            }{%
            \ifthenelse{\equal{\diviseur}{4}}{%
                % Cas 4
                \textbf{Critère de divisibilité par 4 :} Un nombre est divisible par 4 si le nombre formé par ses deux derniers chiffres est divisible par 4.
                
                \textbf{Application à \nombre :}
                \ifthenelse{\nombre < 100}{%
                    Comme \nombre{} est inférieur à 100, nous vérifions directement si \nombre{} est divisible par 4.
                    
                    \ifthenelse{\equal{\intcalcMod{\nombre}{4}}{0}}{%
                        Le reste de la division de \nombre{} par 4 est 0, donc \nombre{} est divisible par 4.
                        }{%
                        Le reste de la division de \nombre{} par 4 est \intcalcMod{\nombre}{4}, donc \nombre{} n'est pas divisible par 4.
                    }%
                    }{%
                    Les deux derniers chiffres de \nombre{} forment le nombre \deuxdernierchiffres.
                    
                    \ifthenelse{\equal{\intcalcMod{\deuxdernierchiffres}{4}}{0}}{%
                        Comme \deuxdernierchiffres{} est divisible par 4, \nombre{} est divisible par 4.
                        }{%
                        Comme \deuxdernierchiffres{} n'est pas divisible par 4 (reste \intcalcMod{\deuxdernierchiffres}{4}), \nombre{} n'est pas divisible par 4.
                    }%
                }%
                }{%
                \ifthenelse{\equal{\diviseur}{5}}{%
                    % Cas 5
                    \textbf{Critère de divisibilité par 5 :} Un nombre est divisible par 5 si son dernier chiffre est 0 ou 5.
                    
                    \textbf{Application à \nombre :} Le dernier chiffre de \nombre{} est \dernierchiffre.
                    
                    \ifthenelse{\equal{\intcalcMod{\nombre}{5}}{0}}{%
                        Comme le dernier chiffre est \dernierchiffre{} (qui est
                        \ifthenelse{\equal{\dernierchiffre}{0}}{0}{5}), \nombre{} est divisible par 5.
                        }{%
                        Comme le dernier chiffre est \dernierchiffre, qui n'est ni 0 ni 5, \nombre{} n'est pas divisible par 5.
                    }%
                    }{%
                    \ifthenelse{\equal{\diviseur}{6}}{%
                        % Cas 6
                        \textbf{Critère de divisibilité par 6 :} Un nombre est divisible par 6 s'il est divisible à la fois par 2 et par 3.
                        
                        \textbf{Application à \nombre :}
                        
                        \textbf{Divisibilité par 2 :} Un nombre est divisible par 2 si son dernier chiffre est pair.
                        
                        \ifthenelse{\equal{\intcalcMod{\nombre}{2}}{0}}{%
                            Le dernier chiffre de \nombre{} est \dernierchiffre, qui est pair, donc \nombre{} est divisible par 2.
                            }{%
                            Le dernier chiffre de \nombre{} est \dernierchiffre, qui n'est pas pair, donc \nombre{} n'est pas divisible par 2.
                        }%
                        
                        \textbf{Divisibilité par 3 :} Un nombre est divisible par 3 si la somme de ses chiffres est divisible par 3.
                        
                        \ifthenelse{\equal{\intcalcMod{\nombre}{3}}{0}}{%
                            La somme des chiffres de \nombre{} est divisible par 3, donc \nombre{} est divisible par 3.
                            }{%
                            La somme des chiffres de \nombre{} n'est pas divisible par 3, donc \nombre{} n'est pas divisible par 3.
                        }%
                        
                        \textbf{Conclusion :}
                        
                        \ifthenelse{\equal{\intcalcMod{\nombre}{6}}{0}}{%
                            Comme \nombre{} est divisible à la fois par 2 et par 3, \nombre{} est divisible par 6.
                            }{%
                            \ifthenelse{\equal{\intcalcMod{\nombre}{2}}{0}}{%
                                Comme \nombre{} n'est pas divisible par 3, \nombre{} n'est pas divisible par 6.
                                }{%
                                Comme \nombre{} n'est pas divisible par 2, \nombre{} n'est pas divisible par 6.
                            }%
                        }%
                        }{%
                        \ifthenelse{\equal{\diviseur}{9}}{%
                            % Cas 9
                            \textbf{Critère de divisibilité par 9 :} Un nombre est divisible par 9 si la somme de ses chiffres est divisible par 9.
                            
                            \textbf{Application à \nombre :}
                            
                            Pour vérifier si \nombre{} est divisible par 9, calculons la somme de ses chiffres.
                            
                            % Pour les besoins de l'exemple, nous affichons directement le résultat
                            \ifthenelse{\equal{\intcalcMod{\nombre}{9}}{0}}{%
                                La somme des chiffres de \nombre{} est divisible par 9, donc \nombre{} est divisible par 9.
                                }{%
                                La somme des chiffres de \nombre{} n'est pas divisible par 9, donc \nombre{} n'est pas divisible par 9.
                            }%
                            }{%
                            % Diviseur invalide
                            %\textbf{Erreur :} Le diviseur \diviseur{} n'est pas dans la liste autorisée (2, 3, 4, 5, 6, 9).
                            Aucun critère ne fonctionne.
                        }%
                    }%
                }%
            }%
        }%
    }%
}


\ExplSyntaxOn

% Définition de la commande \para
% Si \na > 0, retourne \na ; si \na < 0, retourne (\na)
\cs_new:Npn \para #1
{
    \int_compare:nNnTF {#1} > {0}
    {#1} % Si #1 est positif, retourne #1
    {(#1)} % Si #1 est négatif, retourne (#1)
}

% Définition de la commande \expa
% Si \puisxa = 0, retourne \para x^{\puisxa}, sinon retourne \para
\cs_new:Npn \expa #1 #2
{
    \int_compare:nNnTF {#2} = {0}
    {\para{#1}} % Si la puissance est 0, affiche avec x^{#2}
    {
        \int_compare:nNnTF {#2} = {1}
        {\para{#1}~x}
        {\para{#1}~x^{#2}}
    } % Sinon, retourne simplement \para
}
\cs_new:Npn \expan #1 #2
{
    \int_compare:nNnTF {#2} = {0}
    {#1} % Si la puissance est 0, affiche avec x^{#2}
    {
        \int_compare:nNnTF {#2} = {1}
        {#1~x}
        {#1~x^{#2}}
    } % Sinon, retourne simplement \para
}
\cs_new:Npn \finalsigne #1
{
    \int_compare:nNnTF {#1} > {0}
    {+} % Si la puissance est 0, affiche avec x^{#2}
    {-} % Sinon, retourne simplement \para
}
\ExplSyntaxOff
% Définition de la couleur personnalisée
\definecolor{customBackground}{RGB}{0, 43, 54}

\def\boiteQFdetcouleur{customBackground}%{green!75!black}

\newcommand\boiteQFdet[2]{
    \begin{tcolorbox}[nobeforeafter,title=#1,halign title=flush left,fonttitle=\bfseries,colbacktitle=\boiteQFdetcouleur,colframe=customBackground,coltitle=white,colback=white,left=0.2pt,right=0.2pt,width=17.5cm]
        #2
    \end{tcolorbox}
}
\newcommand{\calculatriceAutorisee}{}
% On définit un booléen pour la détection
\newif\ifcalcAllowed
\newcommand\boiteQFQ[2]{%
    % Vérifie si #2 contient la chaîne "\calculatriceAutorisee"
    \IfSubStr{\detokenize{#2}}{\string\calculatriceAutorisee}{%
        % Si trouvé : titre avec l'icône
        \calcAllowedtrue
        }{%
        \calcAllowedfalse
    }%
    %\def\mytitle{#1}
    \begin{tcolorbox}[%
        nobeforeafter,
        title=#1,
        halign title=flush left,
        fonttitle=\bfseries,
        colbacktitle=customBackground,
        colframe=customBackground,
        coltitle=white,
        colback=white,
        left=0.2pt,
        right=0.2pt,
        height=4cm,
        width=8.2cm,
        % L'overlay ne modifie pas la géométrie du titre
        overlay={%
            \ifcalcAllowed
            % Place l'image en overlay à la droite du titre (ajustez xshift/yshift si besoin)
        \node[anchor=base west] at ([xshift=2.5cm,yshift=-0.325cm]title.west) {%xshift=-0.93cm,yshift=-0.325cm]title.east
            \includegraphics[width=0.8cm,height=0.8cm]{C:/Users/Utilisateur/Desktop/Faire/Macros/Programmes_de_facilitation/QF_generator/images/je_calcule.png}%
        };
        \fi
    }%
]
%\begin{itembox}\textbf{#1}\end{itembox}
#2
\end{tcolorbox}
}
%\def\boitereponsecouleur{gray!65!red}
\def\boitereponsecouleur{red!75!black}
\newcommand\boiteQFA[1]{
\begin{tcolorbox}[nobeforeafter,title=Réponses :,halign title=center,fonttitle=\bfseries,colbacktitle=\boitereponsecouleur,colframe=customBackground,coltitle=white,colback=white,left=5pt,right=0.2pt,height=4cm,width=8.2cm]
#1
\end{tcolorbox}
}

\def\myiconpath{}

\getRandomIconPath

\renewcommand\boiteQFA[1]{
\begin{tcolorbox}[nobeforeafter, title=Réponses :, halign title= flush left,
fonttitle=\bfseries, colbacktitle=\boitereponsecouleur,
coltitle=white, colback=white,colframe=customBackground, left=5pt, right=0.2pt,
height=4cm, width=8.2cm, enhanced, overlay={
    \node[anchor=center] at (frame.north west) [xshift=-0.5cm, yshift=0cm]
    {
        %\includegraphics[width=1.5cm]{../../../images/sticker_eclair.png}
        \includegraphics[width=1.5cm]{\myiconpath }
        %\randomincludegraphics[2cm][0cm][0cm]
    };
}]
#1
\end{tcolorbox}
}
\newcommand\boiteQzeroFA[1]{
\begin{tcolorbox}[nobeforeafter, title=Réponses :, halign title= flush left,
fonttitle=\bfseries, colbacktitle=\boitereponsecouleur,
coltitle=white, colback=white,colframe=customBackground, left=5pt, right=0.2pt,
height=4cm, width=8.2cm, enhanced, overlay={
\node[anchor=center] at (frame.north west) [xshift=-0.5cm, yshift=0cm]
{
    \includegraphics[width=1.5cm]{C:/Users/Utilisateur/Desktop/Faire/Macros/Programmes_de_facilitation/QF_generator/images/sticker_eclair.png}%{../../../images/sticker_eclair.png}
    %\includegraphics[width=1.5cm]{\myiconpath }
    %\randomincludegraphics[2cm][0cm][0cm]
};
}]
\begin{center}
\includegraphics[height=2.8cm]{\myiconpath }%
\end{center}
\end{tcolorbox}
}
\newtcolorbox{itembox}[1][]{
baseline,
nobeforeafter, % Supprime les espaces avant et après
colframe=customBackground,
colback=customBackground,
coltext=white,
coltitle=white,
boxsep=1pt,
boxrule=0.5pt,
arc=2pt,
left=2pt,
right=2pt,
top=1pt,
bottom=1pt,
width=0.7cm,
#1
}
\usepackage{graphicx}
\usepackage{fontawesome5}
\usepackage{siunitx} % Pour la notation avec virgule
\usepackage{xfp}
% Configurer siunitx pour utiliser la virgule comme séparateur décimal
%\sisetup{output-decimal-marker={,}}
\newcommand{\cperso}[1]{\num{\fpeval{ #1 }}}
% Configuration locale pour le formatage des nombres
\sisetup{%
%  round-mode=places,
%  round-precision=5,
output-decimal-marker={,},
%  group-separator={}
}

% Définir la commande \response
\newcommand{\fracrep}[2]{%
\FPeval{\result}{round(#1/#2:8)} % Calculer la fraction et arrondir à 4 décimales
\num{\result} % Afficher le résultat avec une virgule
}
\opset{decimalsepsymbol={,}}
%\newcommand{\degree}{^\circ}
\begin{document}
%niveau estimé : 1ere
%date de projection : 24_09_2025
%themes abordés : Polynomes degré 2, Polynomes degré 2, Puissances
\pagecolor{customBackground}
\vspace{-0.5cm}
\begin{multicols}{2}
\boiteQFQ{Question 1 :}{
%version 7 - Court

\def\a{7}
\def\b{5}
\def\typeIR{3}
\def\positionx{2}


\acc{Factoriser} l'expression suivante en utilisant une identité remarquable :
\begin{center}$
\IfEq{\typeIR}{1}{
    \IfEq{\positionx}{1}{
        \affiche{\fpeval{\a^2}}x^2 + \affiche{\fpeval{2*\a*\b}}x + \affiche{\fpeval{\b^2}}
        }{
        \affiche{\fpeval{\a^2}} + \affiche{\fpeval{2*\a*\b}}x + \affiche{\fpeval{\b^2}}x^2
    }
    }{
    \IfEq{\typeIR}{2}{
        \IfEq{\positionx}{1}{
            \affiche{\fpeval{\a^2}}x^2 - \affiche{\fpeval{2*\a*\b}}x + \affiche{\fpeval{\b^2}}
            }{
            \affiche{\fpeval{\a^2}} - \affiche{\fpeval{2*\a*\b}}x + \affiche{\fpeval{\b^2}}x^2
        }
        }{
        \IfEq{\positionx}{1}{
            \affiche{\fpeval{\a^2}}x^2 - \affiche{\fpeval{\b^2}}
            }{
            \affiche{\fpeval{\a^2}} - \affiche{\fpeval{\b^2}}x^2
        }
    }
}$\end{center}
}
\boiteQFQ{Question 2 :}{
%version 4 - Court

\def\coeffa{4}
\def\coeffb{-14}
\def\coeffc{11}
\def\mode{1}


\def\myalpha{\fpeval{-\coeffb/(2*\coeffa)}}
\def\mybeta{\fpeval{(\coeffb*\coeffb-4*\coeffa*\coeffc)/(4*\coeffa)}}
\xdef\myalphafrac{\dfrac{\fpeval{-\coeffb}}{\fpeval{2*\coeffa}}}
\xdef\mybetafrac{\dfrac{\fpeval{\coeffb*\coeffb-4*\coeffa*\coeffc}}{\fpeval{4*\coeffa}}}
\xdef\myalphasimp{\displaystyle{\Simplification{\fpeval{-\coeffb}}{\fpeval{2*\coeffa}}}}
\xdef\mybetasimp{\displaystyle{\Simplification{\fpeval{\coeffb*\coeffb-4*\coeffa*\coeffc}}{\fpeval{4*\coeffa}}}}



\ifthenelse{\mode=1}{\acc{Développer} \begin{center}$\affiche{\coeffa}\left(x - \myalphasimp\right)^2 + \mybetasimp$\end{center}}{\acc{Déterminer la forme canonique de} \begin{center}$\affiche{\coeffa}x^2 + \affiche{\coeffb}x + \affiche{\coeffc}$\end{center}}
}
\end{multicols}
\vspace{-0.85cm}
\begin{multicols}{2}

\boiteQFQ{Question 3 :}{
%version 5 - Court

\def\a{9}
\def\b{9}
\def\m{-1}
\def\n{-9}
\def\mode{3}



\textbf{\'Ecrire} sous la forme d'une \acc{puissance d'un nombre} :
\begin{center}
    \ifthenelse{\mode=1}{$\affiche{\a}^{\affiche{\m}} \times \affiche{\a}^{\affiche{\n}} = \dots$}{\ifthenelse{\mode=2}{$\dfrac{\affiche{\a}^{\affiche{\m}}}{\affiche{\a}^{\affiche{\n}}} = \dots$}{\ifthenelse{\mode=3}{$(\affiche{\a}^{\affiche{\m}})^{\affiche{\n}} = \dots$}{\ifthenelse{\mode=4}{$(\affiche{\a} \times \affiche{\b})^{\affiche{\n}} = \dots$}{$\affiche{\a}^{-\affiche{\n}} = \dots$}}}}
\end{center}
}
\boiteQzeroFA{}
\end{multicols}

\newpage
\vspace{-0.5cm}
\begin{multicols}{2}
\boiteQFQ{Question 1 :}{
%version 7 - Court

\def\a{7}
\def\b{5}
\def\typeIR{3}
\def\positionx{2}


\acc{Factoriser} l'expression suivante en utilisant une identité remarquable :
\begin{center}$
    \IfEq{\typeIR}{1}{
        \IfEq{\positionx}{1}{
            \affiche{\fpeval{\a^2}}x^2 + \affiche{\fpeval{2*\a*\b}}x + \affiche{\fpeval{\b^2}}
            }{
            \affiche{\fpeval{\a^2}} + \affiche{\fpeval{2*\a*\b}}x + \affiche{\fpeval{\b^2}}x^2
        }
        }{
        \IfEq{\typeIR}{2}{
            \IfEq{\positionx}{1}{
                \affiche{\fpeval{\a^2}}x^2 - \affiche{\fpeval{2*\a*\b}}x + \affiche{\fpeval{\b^2}}
                }{
                \affiche{\fpeval{\a^2}} - \affiche{\fpeval{2*\a*\b}}x + \affiche{\fpeval{\b^2}}x^2
            }
            }{
            \IfEq{\positionx}{1}{
                \affiche{\fpeval{\a^2}}x^2 - \affiche{\fpeval{\b^2}}
                }{
                \affiche{\fpeval{\a^2}} - \affiche{\fpeval{\b^2}}x^2
            }
        }
    }$\end{center}
}
\boiteQFQ{Question 2 :}{
    %version 4 - Court
    
    \def\coeffa{4}
    \def\coeffb{-14}
    \def\coeffc{11}
    \def\mode{1}
    
    
    \def\myalpha{\fpeval{-\coeffb/(2*\coeffa)}}
    \def\mybeta{\fpeval{(\coeffb*\coeffb-4*\coeffa*\coeffc)/(4*\coeffa)}}
    \xdef\myalphafrac{\dfrac{\fpeval{-\coeffb}}{\fpeval{2*\coeffa}}}
    \xdef\mybetafrac{\dfrac{\fpeval{\coeffb*\coeffb-4*\coeffa*\coeffc}}{\fpeval{4*\coeffa}}}
    \xdef\myalphasimp{\displaystyle{\Simplification{\fpeval{-\coeffb}}{\fpeval{2*\coeffa}}}}
    \xdef\mybetasimp{\displaystyle{\Simplification{\fpeval{\coeffb*\coeffb-4*\coeffa*\coeffc}}{\fpeval{4*\coeffa}}}}
    
    
    
    \ifthenelse{\mode=1}{\acc{Développer} \begin{center}$\affiche{\coeffa}\left(x - \myalphasimp\right)^2 + \mybetasimp$\end{center}}{\acc{Déterminer la forme canonique de} \begin{center}$\affiche{\coeffa}x^2 + \affiche{\coeffb}x + \affiche{\coeffc}$\end{center}}
}

\end{multicols}
\vspace{-0.85cm}
\begin{multicols}{2}

\boiteQFQ{Question 3 :}{
    %version 5 - Court
    
    \def\a{9}
    \def\b{9}
    \def\m{-1}
    \def\n{-9}
    \def\mode{3}
    
    
    
    \textbf{\'Ecrire} sous la forme d'une \acc{puissance d'un nombre} :
    \begin{center}
        \ifthenelse{\mode=1}{$\affiche{\a}^{\affiche{\m}} \times \affiche{\a}^{\affiche{\n}} = \dots$}{\ifthenelse{\mode=2}{$\dfrac{\affiche{\a}^{\affiche{\m}}}{\affiche{\a}^{\affiche{\n}}} = \dots$}{\ifthenelse{\mode=3}{$(\affiche{\a}^{\affiche{\m}})^{\affiche{\n}} = \dots$}{\ifthenelse{\mode=4}{$(\affiche{\a} \times \affiche{\b})^{\affiche{\n}} = \dots$}{$\affiche{\a}^{-\affiche{\n}} = \dots$}}}}
    \end{center}
}
\boiteQFA{
    \vspace{-0.35cm}
    \begin{itemize}[itemsep=0.2em]
        \item[\raisebox{-0.1cm}{\begin{itembox} \textbf{1.} \end{itembox}}] \def\a{7}\def\b{5}\def\typeIR{3}\def\positionx{2}\IfEq{\typeIR}{1}{    \IfEq{\positionx}{1}{        $(\affiche{\a}x + \affiche{\b})^2$    }{        $(\affiche{\a} + \affiche{\b}x)^2$    }}{    \IfEq{\typeIR}{2}{        \IfEq{\positionx}{1}{            $(\affiche{\a}x - \affiche{\b})^2$        }{            $(\affiche{\a} - \affiche{\b}x)^2$        }    }{        \IfEq{\positionx}{1}{            $(\affiche{\a}x + \affiche{\b})(\affiche{\a}x - \affiche{\b})$        }{            $(\affiche{\a} + \affiche{\b}x)(\affiche{\a} - \affiche{\b}x)$        }    }}
    \end{itemize}
}
\end{multicols}
\newpage
\vspace{-0.5cm}
\begin{multicols}{2}
\boiteQFQ{Question 1 :}{
    %version 7 - Court
    
    \def\a{7}
    \def\b{5}
    \def\typeIR{3}
    \def\positionx{2}
    
    
    \acc{Factoriser} l'expression suivante en utilisant une identité remarquable :
    \begin{center}$
        \IfEq{\typeIR}{1}{
            \IfEq{\positionx}{1}{
                \affiche{\fpeval{\a^2}}x^2 + \affiche{\fpeval{2*\a*\b}}x + \affiche{\fpeval{\b^2}}
                }{
                \affiche{\fpeval{\a^2}} + \affiche{\fpeval{2*\a*\b}}x + \affiche{\fpeval{\b^2}}x^2
            }
            }{
            \IfEq{\typeIR}{2}{
                \IfEq{\positionx}{1}{
                    \affiche{\fpeval{\a^2}}x^2 - \affiche{\fpeval{2*\a*\b}}x + \affiche{\fpeval{\b^2}}
                    }{
                    \affiche{\fpeval{\a^2}} - \affiche{\fpeval{2*\a*\b}}x + \affiche{\fpeval{\b^2}}x^2
                }
                }{
                \IfEq{\positionx}{1}{
                    \affiche{\fpeval{\a^2}}x^2 - \affiche{\fpeval{\b^2}}
                    }{
                    \affiche{\fpeval{\a^2}} - \affiche{\fpeval{\b^2}}x^2
                }
            }
        }$\end{center}
    }
    \boiteQFQ{Question 2 :}{
        %version 4 - Court
        
        \def\coeffa{4}
        \def\coeffb{-14}
        \def\coeffc{11}
        \def\mode{1}
        
        
        \def\myalpha{\fpeval{-\coeffb/(2*\coeffa)}}
        \def\mybeta{\fpeval{(\coeffb*\coeffb-4*\coeffa*\coeffc)/(4*\coeffa)}}
        \xdef\myalphafrac{\dfrac{\fpeval{-\coeffb}}{\fpeval{2*\coeffa}}}
        \xdef\mybetafrac{\dfrac{\fpeval{\coeffb*\coeffb-4*\coeffa*\coeffc}}{\fpeval{4*\coeffa}}}
        \xdef\myalphasimp{\displaystyle{\Simplification{\fpeval{-\coeffb}}{\fpeval{2*\coeffa}}}}
        \xdef\mybetasimp{\displaystyle{\Simplification{\fpeval{\coeffb*\coeffb-4*\coeffa*\coeffc}}{\fpeval{4*\coeffa}}}}
        
        
        
        \ifthenelse{\mode=1}{\acc{Développer} \begin{center}$\affiche{\coeffa}\left(x - \myalphasimp\right)^2 + \mybetasimp$\end{center}}{\acc{Déterminer la forme canonique de} \begin{center}$\affiche{\coeffa}x^2 + \affiche{\coeffb}x + \affiche{\coeffc}$\end{center}}
    }
    
\end{multicols}
\vspace{-0.85cm}
\begin{multicols}{2}
    
    \boiteQFQ{Question 3 :}{
        %version 5 - Court
        
        \def\a{9}
        \def\b{9}
        \def\m{-1}
        \def\n{-9}
        \def\mode{3}
        
        
        
        \textbf{\'Ecrire} sous la forme d'une \acc{puissance d'un nombre} :
        \begin{center}
            \ifthenelse{\mode=1}{$\affiche{\a}^{\affiche{\m}} \times \affiche{\a}^{\affiche{\n}} = \dots$}{\ifthenelse{\mode=2}{$\dfrac{\affiche{\a}^{\affiche{\m}}}{\affiche{\a}^{\affiche{\n}}} = \dots$}{\ifthenelse{\mode=3}{$(\affiche{\a}^{\affiche{\m}})^{\affiche{\n}} = \dots$}{\ifthenelse{\mode=4}{$(\affiche{\a} \times \affiche{\b})^{\affiche{\n}} = \dots$}{$\affiche{\a}^{-\affiche{\n}} = \dots$}}}}
        \end{center}
    }
    \boiteQFA{
        \vspace{-0.35cm}
        \begin{itemize}[itemsep=0.2em]
            \item[\raisebox{-0.1cm}{\begin{itembox} \textbf{1.} \end{itembox}}] \def\a{7}\def\b{5}\def\typeIR{3}\def\positionx{2}\IfEq{\typeIR}{1}{    \IfEq{\positionx}{1}{        $(\affiche{\a}x + \affiche{\b})^2$    }{        $(\affiche{\a} + \affiche{\b}x)^2$    }}{    \IfEq{\typeIR}{2}{        \IfEq{\positionx}{1}{            $(\affiche{\a}x - \affiche{\b})^2$        }{            $(\affiche{\a} - \affiche{\b}x)^2$        }    }{        \IfEq{\positionx}{1}{            $(\affiche{\a}x + \affiche{\b})(\affiche{\a}x - \affiche{\b})$        }{            $(\affiche{\a} + \affiche{\b}x)(\affiche{\a} - \affiche{\b}x)$        }    }}
            \item[\raisebox{-0.1cm}{\begin{itembox} \textbf{2.} \end{itembox}}] \def\coeffa{4}\def\coeffb{-14}\def\coeffc{11}\def\mode{1}\def\myalpha{\fpeval{-\coeffb/(2*\coeffa)}}\def\mybeta{\fpeval{(\coeffb*\coeffb-4*\coeffa*\coeffc)/(4*\coeffa)}}\xdef\myalphafrac{\dfrac{\fpeval{-\coeffb}}{\fpeval{2*\coeffa}}}\xdef\mybetafrac{\dfrac{\fpeval{\coeffb*\coeffb-4*\coeffa*\coeffc}}{\fpeval{4*\coeffa}}}\xdef\myalphasimp{\displaystyle{\Simplification{\fpeval{-\coeffb}}{\fpeval{2*\coeffa}}}}\xdef\mybetasimp{\displaystyle{\Simplification{\fpeval{\coeffb*\coeffb-4*\coeffa*\coeffc}}{\fpeval{4*\coeffa}}}}\ifthenelse{\mode=1}{$\affiche{\coeffa}x^2 + \affiche{\coeffb}x + \affiche{\coeffc}$}{$\affiche{\coeffa}\left(x - \myalphasimp\right)^2 + \mybetasimp$}
        \end{itemize}
    }
\end{multicols}
\newpage

\begin{multicols}{2}
    \boiteQFQ{Question 1 :}{
        %version 7 - Court
        
        \def\a{7}
        \def\b{5}
        \def\typeIR{3}
        \def\positionx{2}
        
        
        \acc{Factoriser} l'expression suivante en utilisant une identité remarquable :
        \begin{center}$
            \IfEq{\typeIR}{1}{
                \IfEq{\positionx}{1}{
                    \affiche{\fpeval{\a^2}}x^2 + \affiche{\fpeval{2*\a*\b}}x + \affiche{\fpeval{\b^2}}
                    }{
                    \affiche{\fpeval{\a^2}} + \affiche{\fpeval{2*\a*\b}}x + \affiche{\fpeval{\b^2}}x^2
                }
                }{
                \IfEq{\typeIR}{2}{
                    \IfEq{\positionx}{1}{
                        \affiche{\fpeval{\a^2}}x^2 - \affiche{\fpeval{2*\a*\b}}x + \affiche{\fpeval{\b^2}}
                        }{
                        \affiche{\fpeval{\a^2}} - \affiche{\fpeval{2*\a*\b}}x + \affiche{\fpeval{\b^2}}x^2
                    }
                    }{
                    \IfEq{\positionx}{1}{
                        \affiche{\fpeval{\a^2}}x^2 - \affiche{\fpeval{\b^2}}
                        }{
                        \affiche{\fpeval{\a^2}} - \affiche{\fpeval{\b^2}}x^2
                    }
                }
            }$\end{center}
        }
        \boiteQFQ{Question 2 :}{
            %version 4 - Court
            
            \def\coeffa{4}
            \def\coeffb{-14}
            \def\coeffc{11}
            \def\mode{1}
            
            
            \def\myalpha{\fpeval{-\coeffb/(2*\coeffa)}}
            \def\mybeta{\fpeval{(\coeffb*\coeffb-4*\coeffa*\coeffc)/(4*\coeffa)}}
            \xdef\myalphafrac{\dfrac{\fpeval{-\coeffb}}{\fpeval{2*\coeffa}}}
            \xdef\mybetafrac{\dfrac{\fpeval{\coeffb*\coeffb-4*\coeffa*\coeffc}}{\fpeval{4*\coeffa}}}
            \xdef\myalphasimp{\displaystyle{\Simplification{\fpeval{-\coeffb}}{\fpeval{2*\coeffa}}}}
            \xdef\mybetasimp{\displaystyle{\Simplification{\fpeval{\coeffb*\coeffb-4*\coeffa*\coeffc}}{\fpeval{4*\coeffa}}}}
            
            
            
            \ifthenelse{\mode=1}{\acc{Développer} \begin{center}$\affiche{\coeffa}\left(x - \myalphasimp\right)^2 + \mybetasimp$\end{center}}{\acc{Déterminer la forme canonique de} \begin{center}$\affiche{\coeffa}x^2 + \affiche{\coeffb}x + \affiche{\coeffc}$\end{center}}
        }
        
    \end{multicols}
    \vspace{-0.85cm}
    \begin{multicols}{2}
        
        \boiteQFQ{Question 3 :}{
            %version 5 - Court
            
            \def\a{9}
            \def\b{9}
            \def\m{-1}
            \def\n{-9}
            \def\mode{3}
            
            
            
            \textbf{\'Ecrire} sous la forme d'une \acc{puissance d'un nombre} :
            \begin{center}
                \ifthenelse{\mode=1}{$\affiche{\a}^{\affiche{\m}} \times \affiche{\a}^{\affiche{\n}} = \dots$}{\ifthenelse{\mode=2}{$\dfrac{\affiche{\a}^{\affiche{\m}}}{\affiche{\a}^{\affiche{\n}}} = \dots$}{\ifthenelse{\mode=3}{$(\affiche{\a}^{\affiche{\m}})^{\affiche{\n}} = \dots$}{\ifthenelse{\mode=4}{$(\affiche{\a} \times \affiche{\b})^{\affiche{\n}} = \dots$}{$\affiche{\a}^{-\affiche{\n}} = \dots$}}}}
            \end{center}
        }
        \boiteQFA{
            \vspace{-0.35cm}
            \begin{itemize}[itemsep=0.2em]
                \item[\raisebox{-0.1cm}{\begin{itembox} \textbf{1.} \end{itembox}}] \def\a{7}\def\b{5}\def\typeIR{3}\def\positionx{2}\IfEq{\typeIR}{1}{    \IfEq{\positionx}{1}{        $(\affiche{\a}x + \affiche{\b})^2$    }{        $(\affiche{\a} + \affiche{\b}x)^2$    }}{    \IfEq{\typeIR}{2}{        \IfEq{\positionx}{1}{            $(\affiche{\a}x - \affiche{\b})^2$        }{            $(\affiche{\a} - \affiche{\b}x)^2$        }    }{        \IfEq{\positionx}{1}{            $(\affiche{\a}x + \affiche{\b})(\affiche{\a}x - \affiche{\b})$        }{            $(\affiche{\a} + \affiche{\b}x)(\affiche{\a} - \affiche{\b}x)$        }    }}
                \item[\raisebox{-0.1cm}{\begin{itembox} \textbf{2.} \end{itembox}}] \def\coeffa{4}\def\coeffb{-14}\def\coeffc{11}\def\mode{1}\def\myalpha{\fpeval{-\coeffb/(2*\coeffa)}}\def\mybeta{\fpeval{(\coeffb*\coeffb-4*\coeffa*\coeffc)/(4*\coeffa)}}\xdef\myalphafrac{\dfrac{\fpeval{-\coeffb}}{\fpeval{2*\coeffa}}}\xdef\mybetafrac{\dfrac{\fpeval{\coeffb*\coeffb-4*\coeffa*\coeffc}}{\fpeval{4*\coeffa}}}\xdef\myalphasimp{\displaystyle{\Simplification{\fpeval{-\coeffb}}{\fpeval{2*\coeffa}}}}\xdef\mybetasimp{\displaystyle{\Simplification{\fpeval{\coeffb*\coeffb-4*\coeffa*\coeffc}}{\fpeval{4*\coeffa}}}}\ifthenelse{\mode=1}{$\affiche{\coeffa}x^2 + \affiche{\coeffb}x + \affiche{\coeffc}$}{$\affiche{\coeffa}\left(x - \myalphasimp\right)^2 + \mybetasimp$}
                \item[\raisebox{-0.1cm}{\begin{itembox} \textbf{3.} \end{itembox}}] \def\a{9}\def\b{9}\def\m{-1}\def\n{-9}\def\mode{3}\ifthenelse{\mode=1}{$\affiche{\a}^{\affiche{\fpeval{\m + \n}}}$}{\ifthenelse{\mode=2}{$\affiche{\a}^{\affiche{\fpeval{\m - \n}}}$}{\ifthenelse{\mode=3}{$\affiche{\a}^{\affiche{\fpeval{\m * \n}}}$}{\ifthenelse{\mode=4}{$\affiche{\a}^{\affiche{\n}} \times \affiche{\b}^{\affiche{\n}}$}{$\dfrac{1}{\affiche{\a}^{\affiche{\n}}}$}}}}
            \end{itemize}
        }
    \end{multicols}
    
    \newpage
    
    %\begin{multicols}{2}
    \boiteQFdet{Solution détaillée de la question 1 :}{
        
        %version 7 - Court
        
        \def\a{7}
        \def\b{5}
        \def\typeIR{3}
        \def\positionx{2}
        
        
        \acc{Factoriser} l'expression suivante en utilisant une identité remarquable :
        \begin{center}$
            \IfEq{\typeIR}{1}{
                \IfEq{\positionx}{1}{
                    \affiche{\fpeval{\a^2}}x^2 + \affiche{\fpeval{2*\a*\b}}x + \affiche{\fpeval{\b^2}}
                    }{
                    \affiche{\fpeval{\a^2}} + \affiche{\fpeval{2*\a*\b}}x + \affiche{\fpeval{\b^2}}x^2
                }
                }{
                \IfEq{\typeIR}{2}{
                    \IfEq{\positionx}{1}{
                        \affiche{\fpeval{\a^2}}x^2 - \affiche{\fpeval{2*\a*\b}}x + \affiche{\fpeval{\b^2}}
                        }{
                        \affiche{\fpeval{\a^2}} - \affiche{\fpeval{2*\a*\b}}x + \affiche{\fpeval{\b^2}}x^2
                    }
                    }{
                    \IfEq{\positionx}{1}{
                        \affiche{\fpeval{\a^2}}x^2 - \affiche{\fpeval{\b^2}}
                        }{
                        \affiche{\fpeval{\a^2}} - \affiche{\fpeval{\b^2}}x^2
                    }
                }
            }$\end{center}
            
            \tikz{\draw[dashed, line width=1pt] (0,0) -- (\linewidth,0);}
            
            \vspace{-0.25cm}\setlength{\columnseprule}{0.4pt}\begin{multicols}{2}
            
            
            %version 7 - solutions
            
            \def\a{7}
            \def\b{5}
            \def\typeIR{3}
            \def\positionx{2}
            
            
            \IfEq{\typeIR}{1}{
                \textbf{Solution :} On reconnaît l'identité remarquable $(a + b)^2 = a^2 + 2ab + b^2$ avec $a>0$ et $b>0$.\\
                
                \IfEq{\positionx}{1}{
                    Ici, $a^2 = \affiche{\fpeval{\a^2}}x^2$ donc $a =\affiche{\a}x$
                    
                    Et $b^2 = \affiche{\fpeval{\b^2}}$ donc $b = \affiche{\b}$
                    
                    Vérifions : $2ab = 2 \times \affiche{\a}x \times \affiche{\b} = \affiche{\fpeval{2*\a*\b}}x$
                    
                    Donc : \encadrer[defi]{$\affiche{\fpeval{\a^2}}x^2 + \affiche{\fpeval{2*\a*\b}}x + \affiche{\fpeval{\b^2}} = (\affiche{\a}x + \affiche{\b})^2$	}
                    }{
                    Ici, $a^2 = \affiche{\fpeval{\a^2}}$ donc $a = \affiche{\a}$
                    
                    Et $b^2 = \affiche{\fpeval{\b^2}}x^2$ donc $b = \affiche{\b}x$
                    
                    Vérifions : $2ab = 2 \times \affiche{\a} \times \affiche{\b}x = \affiche{\fpeval{2*\a*\b}}x$
                    
                    Donc : \encadrer[defi]{$\affiche{\fpeval{\a^2}} + \affiche{\fpeval{2*\a*\b}}x + \affiche{\fpeval{\b^2}}x^2 = (\affiche{\a} + \affiche{\b}x)^2$		}
                }
                }{
                \IfEq{\typeIR}{2}{
                    \textbf{Solution :} On reconnaît l'identité remarquable $(a - b)^2 = a^2 - 2ab + b^2$ avec $a>0$ et $b>0$.\\
                    
                    \IfEq{\positionx}{1}{
                        Ici, $a^2 = \affiche{\fpeval{\a^2}}x^2$ donc $a = \affiche{\a}x$
                        
                        Et $b^2 = \affiche{\fpeval{\b^2}}$ donc $b = \affiche{\b}$\\
                        
                        Vérifions : $2ab = 2 \times \affiche{\a}x \times \affiche{\b} = \affiche{\fpeval{2*\a*\b}}x$ \\
                        
                        Donc : \encadrer[defi]{$\affiche{\fpeval{\a^2}}x^2 - \affiche{\fpeval{2*\a*\b}}x + \affiche{\fpeval{\b^2}} = (\affiche{\a}x - \affiche{\b})^2$	}
                        }{
                        Ici, $a^2 = \affiche{\fpeval{\a^2}}$ donc $a = \affiche{\a}$
                        
                        Et $b^2 = \affiche{\fpeval{\b^2}}x^2$ donc $b = \affiche{\b}x$\\
                        
                        Vérifions : $2ab = 2 \times \affiche{\a} \times \affiche{\b}x = \affiche{\fpeval{2*\a*\b}}x$ \\
                        
                        Donc : \encadrer[defi]{$\affiche{\fpeval{\a^2}} - \affiche{\fpeval{2*\a*\b}}x + \affiche{\fpeval{\b^2}}x^2 = (\affiche{\a} - \affiche{\b}x)^2$		}
                    }
                    }{
                    \textbf{Solution :} On reconnaît l'identité remarquable $a^2 - b^2 = (a + b)(a - b)$ avec $a>0$ et $b>0$.\\
                    
                    \IfEq{\positionx}{1}{
                        Ici, $a^2 = \affiche{\fpeval{\a^2}}x^2$ donc $a = \affiche{\a}x$
                        
                        Et $b^2 = \affiche{\fpeval{\b^2}}$ donc $b = \affiche{\b}$\\
                        
                        Donc :\encadrer[defi]{$\affiche{\fpeval{\a^2}}x^2 - \affiche{\fpeval{\b^2}} = (\affiche{\a}x + \affiche{\b})(\affiche{\a}x - \affiche{\b})$	}
                        }{
                        Ici, $a^2 = \affiche{\fpeval{\a^2}}$ donc $a = \affiche{\a}$
                        
                        Et $b^2 = \affiche{\fpeval{\b^2}}x^2$ donc $b = \affiche{\b}x$\\
                        
                        Donc : \encadrer[defi]{$\affiche{\fpeval{\a^2}} - \affiche{\fpeval{\b^2}}x^2 = (\affiche{\a} + \affiche{\b}x)(\affiche{\a} - \affiche{\b}x)$}
                    }
                }
            }
            
        \end{multicols}
    }
    
    \newpage
    
    \boiteQFdet{Solution détaillée de la question 2 :}{
        
        %version 4 - Court
        
        \def\coeffa{4}
        \def\coeffb{-14}
        \def\coeffc{11}
        \def\mode{1}
        
        
        \def\myalpha{\fpeval{-\coeffb/(2*\coeffa)}}
        \def\mybeta{\fpeval{(\coeffb*\coeffb-4*\coeffa*\coeffc)/(4*\coeffa)}}
        \xdef\myalphafrac{\dfrac{\fpeval{-\coeffb}}{\fpeval{2*\coeffa}}}
        \xdef\mybetafrac{\dfrac{\fpeval{\coeffb*\coeffb-4*\coeffa*\coeffc}}{\fpeval{4*\coeffa}}}
        \xdef\myalphasimp{\displaystyle{\Simplification{\fpeval{-\coeffb}}{\fpeval{2*\coeffa}}}}
        \xdef\mybetasimp{\displaystyle{\Simplification{\fpeval{\coeffb*\coeffb-4*\coeffa*\coeffc}}{\fpeval{4*\coeffa}}}}
        
        
        
        \ifthenelse{\mode=1}{\acc{Développer} \begin{center}$\affiche{\coeffa}\left(x - \myalphasimp\right)^2 + \mybetasimp$\end{center}}{\acc{Déterminer la forme canonique de} \begin{center}$\affiche{\coeffa}x^2 + \affiche{\coeffb}x + \affiche{\coeffc}$\end{center}}
        
        \tikz{\draw[dashed, line width=1pt] (0,0) -- (\linewidth,0);}
        
        
        \vspace{-0.25cm}\setlength{\columnseprule}{0.4pt}\begin{multicols}{2}
        
        %version 4 - solutions
        
        \def\coeffa{4}
        \def\coeffb{-14}
        \def\coeffc{11}
        \def\mode{1}
        
        
        \def\myalpha{\fpeval{-\coeffb/(2*\coeffa)}}
        \def\mybeta{\fpeval{(\coeffb*\coeffb-4*\coeffa*\coeffc)/(-4*\coeffa)}}
        \def\mybetanum{\fpeval{\coeffb*\coeffb-4*\coeffa*\coeffc}}
        \def\mybetaden{\fpeval{-4*\coeffa}}
        \xdef\myalphafrac{\dfrac{\fpeval{-\coeffb}}{\fpeval{2*\coeffa}}}
        \xdef\mybetafrac{\dfrac{\fpeval{\coeffb*\coeffb-4*\coeffa*\coeffc}}{\fpeval{-4*\coeffa}}}
        \xdef\myalphasimp{\displaystyle{\Simplification{\fpeval{-\coeffb}}{\fpeval{2*\coeffa}}}}
        \xdef\mybetasimp{\displaystyle{\Simplification{\fpeval{\coeffb*\coeffb-4*\coeffa*\coeffc}}{\fpeval{-4*\coeffa}}}}
        
        \ifthenelse{\mode=1}{
            \textbf{Rappel :} $(x - \alpha)^2 = x^2 - 2\alpha x + \alpha^2$
            
            Développons :
            \begin{align*}
                &\affiche{\coeffa}(x - \myalphasimp)^2 + \mybetasimp\\
                &= \affiche{\coeffa}\left[x^2 - 2 \times \myalphasimp \times x + \left(\myalphasimp\right)^2\right] + \mybetasimp\\
                &\approx \affiche{\coeffa}x^2 - \affiche{\fpeval{round(\coeffa*2*\myalpha,0)}}x + \affiche{\fpeval{round(\coeffa*\myalpha^2,2)}} + \affiche{\fpeval{round(\mybetanum/\mybetaden,2)}} \\
                &= \affiche{\coeffa}x^2 + \affiche{\coeffb}x + \affiche{\coeffc}
            \end{align*}
            }{
            \textbf{Détermination de la forme canonique :}
            
            Pour $ax^2 + bx + c$, la forme canonique est $a(x - \alpha)^2 + \beta$ avec :
            \begin{multicols}{2}
                \begin{itemize}
                    \item $\alpha = -\dfrac{b}{2a}$
                    \item $\beta = -\dfrac{b^2 - 4ac}{4a}$
                \end{itemize}
                \columnbreak
                
                Avec $a = \affiche{\coeffa}$,
                
                $b = \affiche{\coeffb}$ et $c = \affiche{\coeffc}$ :
            \end{multicols}
            
            \columnbreak
            
            $\alpha = -\dfrac{\affiche{\coeffb}}{2 \times \affiche{\coeffa}} = \myalphafrac =$ \encadrer[defi]{$\myalphasimp$}
            
            $\beta = -\dfrac{\affiche{\coeffb}^2 - 4 \times \affiche{\coeffa} \times \affiche{\coeffc}}{4 \times \affiche{\coeffa}} = -\dfrac{\affiche{\fpeval{\coeffb*\coeffb}} - \affiche{\fpeval{4*\coeffa*\coeffc}}}{\affiche{\fpeval{4*\coeffa}}}$
            
            $\phantom{\beta } = \mybetafrac = $\encadrer[defi]{$\mybetasimp$}
            
            Donc :
            
            \begin{center}
                \encadrer[defi]{$\affiche{\coeffa}x^2 + \affiche{\coeffb}x + \affiche{\coeffc} = \affiche{\coeffa}\left(x - \myalphasimp\right)^2 + \mybetasimp$}
            \end{center}
        }
        
    \end{multicols}
}

%\columnbreak
\newpage


\boiteQFdet{Solution détaillée de la question 3 :}{
    
    %version 5 - Court
    
    \def\a{9}
    \def\b{9}
    \def\m{-1}
    \def\n{-9}
    \def\mode{3}
    
    
    
    \textbf{\'Ecrire} sous la forme d'une \acc{puissance d'un nombre} :
    \begin{center}
        \ifthenelse{\mode=1}{$\affiche{\a}^{\affiche{\m}} \times \affiche{\a}^{\affiche{\n}} = \dots$}{\ifthenelse{\mode=2}{$\dfrac{\affiche{\a}^{\affiche{\m}}}{\affiche{\a}^{\affiche{\n}}} = \dots$}{\ifthenelse{\mode=3}{$(\affiche{\a}^{\affiche{\m}})^{\affiche{\n}} = \dots$}{\ifthenelse{\mode=4}{$(\affiche{\a} \times \affiche{\b})^{\affiche{\n}} = \dots$}{$\affiche{\a}^{-\affiche{\n}} = \dots$}}}}
    \end{center}
    
    \tikz{\draw[dashed, line width=1pt] (0,0) -- (\linewidth,0);}
    
    \vspace{-0.25cm}\setlength{\columnseprule}{0.4pt}\begin{multicols}{2}
    
    %version 5 - solutions
    
    \def\a{9}
    \def\b{9}
    \def\m{-1}
    \def\n{-9}
    \def\mode{3}
    
    
    
    \ifthenelse{\mode=1}{
        \textbf{Formule :} $a^m \times a^n = a^{m+n}$ avec $a = \affiche{\a}$, $m = \affiche{\m}$ et $n = \affiche{\n}$
        
        $\affiche{\a}^{\affiche{\m}} \times \affiche{\a}^{\affiche{\n}} = \affiche{\a}^{\affiche{\m} + \affiche{\n}} = \affiche{\a}^{\affiche{\fpeval{\m + \n}}}$
        }{
        \ifthenelse{\mode=2}{
            \textbf{Formule :} $\dfrac{a^m}{a^n} = a^{m-n}$ avec $a = \affiche{\a}$, $m = \affiche{\m}$ et $n = \affiche{\n}$
            
            $\dfrac{\affiche{\a}^{\affiche{\m}}}{\affiche{\a}^{\affiche{\n}}} = \affiche{\a}^{\affiche{\m} - \affiche{\n}} = \affiche{\a}^{\affiche{\fpeval{\m - \n}}}$
            }{
            \ifthenelse{\mode=3}{
                \textbf{Formule :} $(a^m)^n = a^{mn}$ avec $a = \affiche{\a}$, $m = \affiche{\m}$ et $n = \affiche{\n}$
                
                $(\affiche{\a}^{\affiche{\m}})^{\affiche{\n}} = \affiche{\a}^{\affiche{\m} \times \affiche{\n}} = \affiche{\a}^{\affiche{\fpeval{\m * \n}}}$
                }{
                \ifthenelse{\mode=4}{
                    \textbf{Formule :} $(ab)^n = a^n b^n$ avec $a = \affiche{\a}$, $b = \affiche{\b}$ et $n = \affiche{\n}$
                    
                    $(\affiche{\a} \times \affiche{\b})^{\affiche{\n}} = \affiche{\a}^{\affiche{\n}} \times \affiche{\b}^{\affiche{\n}}$
                    }{
                    \textbf{Formule :} $a^{-n} = \dfrac{1}{a^n}$ avec $a = \affiche{\a}$ et $n = \affiche{\n}$
                    
                    $\affiche{\a}^{-\affiche{\n}} = \dfrac{1}{\affiche{\a}^{\affiche{\n}}}$
    }}}}
    
\end{multicols}
}
%\end{multicols}
\end{document}
